\begin{resumo}
 Plataformas de dados governamentais já consolidam o histórico de execução de
 políticas públicas, mas raramente contam com infraestrutura para colocar
 modelos de aprendizado de máquina em produção de forma confiável. Este
 trabalho propõe uma arquitetura de MLOps que opera sobre um \textit{Data
 Lakehouse} organizado nas camadas da arquitetura medalhão, em vez de construir
 uma plataforma de aprendizado de máquina paralela. A camada bronze preserva o
 dado bruto como recebido da fonte; a camada prata reúne os dados limpos e
 harmonizados, dos quais derivam o painel histórico e o rótulo; e a camada ouro
 concentra as features calculadas na data de cada observação, os escores do
 modelo e as métricas de monitoramento. A arquitetura é definida por
 responsabilidades, independentes de ferramenta: a camada de transformação
 concentra todo o cálculo sobre dados, o ambiente de modelagem apenas treina,
 avalia e escora, o orquestrador coordena as etapas e o registro de
 experimentos guarda métricas e versões de modelo. Os escores voltam ao
 \textit{lakehouse} pela camada bronze e percorrem as demais camadas como
 qualquer outra fonte, sob os mesmos testes e o mesmo catálogo. Cada treino
 congela seu dataset com um manifesto que permite refazer o resultado a partir
 da versão do código e dos dados. A pesquisa segue a \textit{Design Science
 Research}, e a arquitetura é instanciada na previsão de paralisação de obras
 do Programa Minha Casa, Minha Vida na plataforma Gov Hub Cidades, do
 Ministério das Cidades, tendo como referência de comparação a regra de
 limiares fixos usada hoje pela fiscalização. Nesta primeira etapa, foram
 definidos os requisitos e o projeto da arquitetura e realizado o levantamento
 inicial dos dados, que revelou que as trocas de fonte ao longo da série
 histórica restringem o período em que o rótulo de paralisação é confiável. A
 avaliação, prevista para a segunda etapa, considera a aderência a princípios
 de MLOps da literatura, a prontidão para produção, a reprodutibilidade, a
 ausência de vazamento temporal e o desempenho do modelo frente à regra
 vigente.

 \vspace{\onelineskip}

 \noindent
 \textbf{Palavras-chave}: MLOps. dados governamentais. \textit{data lakehouse}.
 arquitetura medalhão. aprendizado de máquina.
\end{resumo}
